\scanpage{31}\mspage{29}
On se propose d'appliquer \unsure{la|le} \deleted{``calcul'' préc.} \inserted{\unsure{théorie|théorème} du} \S\,2
au ``calcul'' \unsure{des|de}
\[
\operatorname{Hom}_{\unsure{\Sigma|Z}_X}(a_X,b_X),
\]
on a donc une \unsure{bij.} can.
\[
\boxed{\operatorname{Hom}_{\Sigma_X}(a_X,b_X)\simeq\pi_0\underline{\mathrm{Ch}}_{\deleted{\unsure{M(X)}}\Sigma_X}\bigl((\unsure{M|\mathcal{M}}(X);a_X,b_X\bigr).}
\]
On se propose d'expliciter la catégorie des
chemins du second membre, On a, abstraction
faite des données de $\Sigma,a,b$ dans $M$, la
\deleted{$\unsure{\underline{\mathrm{Ch}}_{\Sigma_X}}$}
expression pour la catégorie de \uline{tous} les
chemins dans $M(X)$
\[
\begin{aligned}
\underline{\mathrm{Ch}}(M(X))
&\simeq\underline{\mathrm{Hom}}\bigl(
\underset{\substack{\text{pro-cat.}\\(\unsure{\mathscr{L}|\mathcal{L}}_n)_{n\in\unsure{\mathbb{N}|N}}\ \\ \deleted{\unsure{n\ \mathrm{pair}}}}}{\unsure{\mathbb{L}|\mathscr{L}}},
\underbrace{M(X)}_{\underline{\mathrm{Hom}}(X,M)}\bigr)\quad )\\
&\overset{\mathrm{def}}{\simeq}\varinjlim_n
\underrightarrow{\underline{\mathrm{Hom}}}\bigl(\mathscr{L}_n,\underbrace{\overbrace{M(X)}}\bigr)\\
&\simeq\varinjlim_n\underbrace{\underline{\mathrm{Hom}}(\mathscr{L}_n\times X,M)}_{
\simeq\underline{\mathrm{Hom}}\bigl(X,
\underbrace{\underline{\mathrm{Hom}}(\mathscr{L}_n,M)}_{\substack{\Vert\,\mathrm{def}\\\underline{\mathrm{Ch}}_n(M)}}\bigr)}\\
&\simeq\varinjlim_n\underline{\mathrm{Hom}}\bigl(X,\underline{\mathrm{Ch}}_n(M)\bigr)\\
&\simeq\underline{\mathrm{Hom}}^!\bigl(X,\underline{\mathrm{Ch}}(M)\bigr)
\end{aligned}
\]
où le signe $!$ signifie qu'on prend la sous-catégorie pleine des foncteurs \deleted{$\unsure{X}$}
\[
X\unsure{\mathord{-}|\to}\underline{\mathrm{Ch}}(M)
\]
% archive sheet number: 30
